\section{Adaptation policies}
\er{the policies need better framing: we must present first what is the control/feedback loop and say we will present a few variants. We must describe clearly what is the problem first. We could separate between the data available to the nodes. I insist that having a table linking the currently-running protocol and the available data would be very useful, as it drives the type of adaptation policy we can have.}

\er{the proposed strategies (algorithms) integrate the threshold-based rules directly in the algorithm. It would be clearer to put them separately as predicates provided to the policies, and that can be reused across multiple decision-making strategies (e.g. the simple local one vs the complex collective one).}

\er{when using adaptation policies, it is necessary to use hysteresis to avoid constant oscillations between two protocols. We can also include a grace period where a node keeps its protocol for a time before being allowed to switch again.}
\raziel{proposal of simple adaptation policy}
\begin{enumerate}

\item all nodes are running $A_{i}$

\item source node forward message $W_{msg}$ (wiling to change)
	\begin{itemize}
	\item this action may take place when collisions increase (getting locally from MAC layer) or when there is a peak of energy consumption?
	\end{itemize}
	\begin{enumerate}
		\item start a timer $T_{will}$
		\item start exchange of control messages
	\end{enumerate}

\item reception of message $W_{msg}$
	\begin{enumerate}
		\item forward $W_{msg}$
		\item start sending control messages to eventually compute the local clustering coefficient (LCC) ($c(l)$)
		\item log velocity $V_{0}$
	\end{enumerate}

\item all nodes in the communication area receive $W_{msg}$ based in the way $A_{i}$ works

\item when timer $T_{will}$ (see 2.a) expires at the source node
	\begin{enumerate}
		\item log velocity $V_{1}$
		\item get $n(a(V_{0}, V_{1}))$ where $a$ is the acceleration
		\item get LCC
		\item stop sending beacons
		\item IF $z(l, a) \approx 1.0$ THEN (node is at dense/static zone)
		\begin{enumerate}
			\item create message $C_{msg}=M(z(l, a), A_{cds})$ where $A_{cds}$ refers to an algorithm that relies in an overlay
			\item forward $C_{msg}$
			\item switch to $A_{cds}$
		\end{enumerate}
		\item IF $z(l, a) \approx -1.0$ DO steps e.[i - iii] with a non-overlay based algorithm identifier $A_{f}$ (node is at sparse/dynamic zone)
	\end{enumerate}

\item Reception of $C_{msg}$
	\begin{enumerate}
		\item IF $A_{i} \neq C_{msg}[A] \land z(l,a) - C_{msg}[z] \approx 0$ THEN
		\begin{itemize}
			\item CASE: sender are receiver are within the same zone
			\item forward message ${C_{msg}}$ as it was received
			\item start running new algorithm $C_{msg}[A]$ instead of $A_{i}$
		\end{itemize}
	\end{enumerate}
		\begin{enumerate}
		\item IF $A_{i} \neq C_{msg}[A] \land z(l,a) - C_{msg}[z] < 0$ THEN
		\begin{itemize}
			\item CASE: first discovered border node
			\item create message $C_{msg}=M(z(l, a), A_{flooding})$; \textbf{NOTE}: a node located at the border intents to reach other nodes at $Z_{sparse}$
			\item forward $C_{msg}$
			\item start running new algorithm $A_{flooding}$ which was taken from catalogs of algorithms
			\item \textbf{NOTE:} when a border node forwards $C_{msg}$ its identifier remains the same; meaning that receivers at $Z_{dense}$ ignore such message
		\end{itemize}
	\end{enumerate}

\end{enumerate}